\documentclass[12pt,a4paper]{article}

% ----- Encoding & Language -----
\usepackage[utf8]{inputenc}
\usepackage[T1]{fontenc}

% ----- Mathematics -----
\usepackage{amsmath, amssymb, amsthm}
\usepackage{mathtools}

% ----- Layout -----
\usepackage[a4paper, margin=2.5cm]{geometry}
\usepackage{setspace}
\onehalfspacing

% ----- Colours & Graphics -----
\usepackage[dvipsnames]{xcolor}
\definecolor{mainblue}{RGB}{30,75,182}
\definecolor{darkblue}{RGB}{31,56,100}
\definecolor{lightblue}{RGB}{217,226,243}
\definecolor{codebg}{RGB}{245,247,252}
\definecolor{codefg}{RGB}{31,56,100}

% ----- Headers & Footers -----
\usepackage{fancyhdr}
\pagestyle{fancy}
\fancyhf{}
\fancyhead[L]{\small\color{mainblue}\textit{MDPL26 --- Matemática Discreta, Pós-laboral}}
\fancyhead[R]{\small\color{mainblue}\textit{Ponto 3 --- Alíneas (a), (b), (c)}}
\fancyfoot[C]{\small\color{gray}\thepage}
\renewcommand{\headrulewidth}{0.4pt}

% ----- Section Formatting -----
\usepackage{titlesec}
\titleformat{\section}{\large\bfseries\color{darkblue}}{}{0em}{}[\color{mainblue}\titlerule]
\titleformat{\subsection}{\normalsize\bfseries\color{mainblue}}{}{0em}{}
\titleformat{\subsubsection}{\normalsize\itshape\color{mainblue}}{}{0em}{}

% ----- Code Listings -----
\usepackage{listings}
\lstset{
  backgroundcolor=\color{codebg},
  basicstyle=\ttfamily\small\color{codefg},
  frame=single,
  framerule=0.4pt,
  rulecolor=\color{mainblue!40},
  breaklines=true,
  columns=flexible,
  keepspaces=true,
  xleftmargin=1em,
  xrightmargin=0.5em,
  aboveskip=6pt,
  belowskip=6pt,
}

% ----- Tables -----
\usepackage{booktabs}
\usepackage{array}
\usepackage{colortbl}

% ----- Enumeration -----
\usepackage{enumitem}
\setlist[itemize]{leftmargin=1.5em, itemsep=2pt, topsep=4pt}

% ----- Hyperlinks -----
\usepackage[hidelinks]{hyperref}

% ----- Diagrams -----
\usepackage{tikz-cd}
\usepackage{tikz}

% ----- Title -----
\title{
  \vspace{-1cm}
  {\large Politécnico de Leiria}\\[0.3em]
  {\color{darkblue}\rule{\linewidth}{1.5pt}}\\[0.4em]
  {\LARGE\bfseries\color{darkblue} Matemática Discreta --- MDPL26}\\[0.3em]
  {\large\color{mainblue} Ponto 3 --- Funtores e Transformações Naturais}\\[0.2em]
  {\normalsize\color{gray} Alíneas (a), (b) e (c)}\\[0.3em]
  {\color{darkblue}\rule{\linewidth}{1.5pt}}
}
\author{}
\date{}

% ----- Custom commands -----
\newcommand{\catname}[1]{\mathbf{#1}}
\newcommand{\Fin}{\catname{Fin}}
\newcommand{\Vect}{\catname{Vect}}
\newcommand{\Sub}{\catname{Sub}}
\newcommand{\Par}{\catname{Par}}
\newcommand{\C}{\mathbb{C}}
\newcommand{\N}{\mathbb{N}}
\newcommand{\parto}{\rightharpoonup}  % partial map arrow

% Coloured box for definitions/remarks
\usepackage{tcolorbox}
\tcbuselibrary{skins}
\newtcolorbox{defbox}[1][]{
  colback=lightblue!50,
  colframe=mainblue,
  fonttitle=\bfseries\small,
  title=#1,
  boxrule=0.6pt,
  arc=3pt,
  left=6pt, right=6pt, top=4pt, bottom=4pt,
}

% ============================================================
\begin{document}
\maketitle
\thispagestyle{fancy}

\begin{abstract}
\noindent
Este documento apresenta, para cada par de categorias indicado nas alíneas (a), (b) e (c) do Ponto~3, os funtores covariantes e contravariantes possíveis, bem como as transformações naturais entre eles, com recurso a descrições simplificadas em MATLAB/Octave.
\end{abstract}

\tableofcontents
\vspace{1em}

% ============================================================
\section{Alínea (a) --- \texorpdfstring{$\Fin$}{Fin} e \texorpdfstring{$\Vect$}{Vect}}

\subsection{Descrição das categorias (modelo MATLAB/Octave)}

\begin{defbox}[Categoria $\Fin$]
Um objecto é representado pelo seu cardinal $n \in \N$. Um morfismo $f \colon n \to m$ é um vector de inteiros de comprimento $n$ com entradas em $\{1,\ldots,m\}$, onde $f(i)$ indica a imagem do elemento $i$.
\end{defbox}

\begin{lstlisting}[language=Matlab, caption={Objeto e morfismo em $\Fin$}]
n = 3;          % objecto: conjunto {1,2,3}
f = [2, 2, 1];  % morfismo: f(1)=2, f(2)=2, f(3)=1
\end{lstlisting}

\begin{defbox}[Categoria $\Vect$]
Um objecto é um inteiro $d \geq 0$ representando a dimensão do espaço vectorial sobre $\C$. Um morfismo $L \colon d_1 \to d_2$ é uma matriz complexa $d_2 \times d_1$.
\end{defbox}

\begin{lstlisting}[language=Matlab, caption={Objeto e morfismo em $\Vect$}]
d = 3;              % objecto: C^3
L = [1 0 0; 0 1 0]; % morfismo: aplicacao linear C^3 -> C^2
\end{lstlisting}

% ---
\subsection{Funtor covariante \texorpdfstring{$F \colon \Fin \to \Vect$}{F: Fin -> Vect} (funtor livre)}

O \emph{funtor livre} $F \colon \Fin \to \Vect$ envia cada conjunto finito $n$ ao espaço vectorial $\C^n$ gerado pelos seus elementos como base canónica, e cada aplicação total $f \colon n \to m$ à aplicação linear $F(f) \colon \C^n \to \C^m$ definida na base canónica por
\[
  e_i \;\longmapsto\; e_{f(i)}.
\]
A imagem do funtor num morfismo $f$ é a matriz de permutação/incidência:

\begin{lstlisting}[language=Matlab, caption={Funtor livre $F$ aplicado a um morfismo}]
function L = free_functor(f, m)
  n = length(f);
  L = zeros(m, n);
  for i = 1:n
    L(f(i), i) = 1;
  end
end
\end{lstlisting}

\noindent\textbf{Propriedades do funtor $F$:}
\begin{itemize}
  \item $F(\mathrm{id}_n) = I_n$ (matriz identidade $n \times n$).
  \item $F(g \circ f) = F(g) \cdot F(f)$ (preservação da composição).
  \item $F$ envia aplicações injectivas a aplicações lineares injectivas.
  \item $F$ envia aplicações sobrejectivas a aplicações lineares sobrejectivas.
\end{itemize}

O diagrama de naturalidade toma a forma:
\[
\begin{tikzcd}
  F(n) \arrow[r, "F(f)"] \arrow[d, "\eta_n"'] & F(m) \arrow[d, "\eta_m"] \\
  G(n) \arrow[r, "G(f)"']                     & G(m)
\end{tikzcd}
\]

% ---
\subsection{Funtor de esquecimento \texorpdfstring{$U \colon \Vect \to \Fin$}{U: Vect -> Fin}}

O funtor de esquecimento $U \colon \Vect \to \Fin$ envia cada espaço vectorial de dimensão $d$ ao conjunto finito de índices $\{1,\ldots,d\}$ (a base canónica). Uma aplicação linear $L \colon \C^{d_1} \to \C^{d_2}$ é esquecida enquanto estrutura linear. Note-se que este processo requer a escolha de uma base, pelo que $U$ é bem definido apenas após fixar bases em cada objecto.

% ---
\subsection{Transformações naturais \texorpdfstring{$\eta \colon F \Rightarrow G$}{eta: F => G}}

Dados dois funtores $F, G \colon \Fin \to \Vect$, uma \emph{transformação natural} $\eta \colon F \Rightarrow G$ é uma família de morfismos em $\Vect$,
\[
  \eta_n \colon F(n) \to G(n), \quad \text{para cada objecto } n \in \Fin,
\]
tal que para cada morfismo $f \colon n \to m$ em $\Fin$ o seguinte quadrado comuta:
\[
  G(f) \circ \eta_n \;=\; \eta_m \circ F(f).
\]
Em termos matriciais (MATLAB/Octave), a condição de naturalidade é:

\begin{lstlisting}[language=Matlab, caption={Condição de naturalidade em $\Vect$}]
% G_f, F_f: matrizes de F(f) e G(f); eta_n, eta_m: componentes de eta
norm(G_f * eta_n - eta_m * F_f) < tol  % deve ser verdadeiro
\end{lstlisting}

\noindent\textbf{Exemplo:} Seja $F$ o funtor livre e $G$ um funtor que envia $n$ a $\C^n$ via uma base diferente. Então $\eta_n$ é simplesmente uma matriz de mudança de base invertível em cada componente, compatível com os morfismos.

% ============================================================
\section{Alínea (b) --- \texorpdfstring{$\Fin$}{Fin} e \texorpdfstring{$\Sub$}{Sub}}

\subsection{Descrição das categorias (modelo MATLAB/Octave)}

\begin{defbox}[Categoria $\Sub$]
Um objecto é um vector de inteiros (os elementos do conjunto). Um morfismo $A \to B$ existe apenas se $A \subseteq B$, e é representado pelo vector lógico de inclusão.
\end{defbox}

\begin{lstlisting}[language=Matlab, caption={Morfismo de inclusão em $\Sub$}]
A = [1, 3];  B = [1, 2, 3, 4];
is_sub = all(ismember(A, B));  % verifica se A c= B
incl   = ismember(B, A);       % vector logico da inclusao
\end{lstlisting}

% ---
\subsection{Funtor covariante \texorpdfstring{$P \colon \Fin \to \Sub$}{P: Fin -> Sub} (conjunto das partes)}

O \emph{funtor potência} $P \colon \Fin \to \Sub$ envia cada conjunto finito $n$ ao conjunto das suas partes $\mathcal{P}(n) = \{S \subseteq \{1,\ldots,n\}\}$, e cada aplicação total $f \colon n \to m$ à aplicação
\[
  P(f) \colon \mathcal{P}(n) \to \mathcal{P}(m), \quad S \;\longmapsto\; f(S) = \{f(i) : i \in S\}.
\]

\begin{lstlisting}[language=Matlab, caption={Funtor potência $P$}]
function PS = power_set(n)
  PS = {};
  for k = 0:2^n-1
    PS{end+1} = find(bitand(k, 2.^(0:n-1)));
  end
end

function img = direct_image(f, S)
  img = unique(f(S));  % imagem directa de S por f
end
\end{lstlisting}

\noindent Este funtor preserva inclusões: se $S \subseteq T$ então $f(S) \subseteq f(T)$, pelo que está bem definido em $\Sub$ como codomínio.

% ---
\subsection{Funtor de esquecimento \texorpdfstring{$U \colon \Sub \to \Fin$}{U: Sub -> Fin}}

O funtor de esquecimento $U \colon \Sub \to \Fin$ envia cada objecto de $\Sub$ (um subconjunto $A \subseteq B$) ao conjunto finito $A$ (ignorando a estrutura de inclusão), e cada morfismo (inclusão $A \hookrightarrow B$) à correspondente aplicação total em $\Fin$.

% ---
\subsection{Funtor contravariante \texorpdfstring{$P^* \colon \Fin^{\mathrm{op}} \to \Sub$}{P*: Fin^op -> Sub} (imagem inversa)}

O \emph{funtor imagem inversa} $P^* \colon \Fin^{\mathrm{op}} \to \Sub$ envia cada aplicação $f \colon n \to m$ à aplicação
\[
  P^*(f) \colon \mathcal{P}(m) \to \mathcal{P}(n), \quad T \;\longmapsto\; f^{-1}(T) = \{i \in n : f(i) \in T\}.
\]
Este funtor é \emph{contravariante} porque inverte a direcção dos morfismos.

\begin{lstlisting}[language=Matlab, caption={Funtor contravariante $P^*$ (pré-imagem)}]
function preimg = inverse_image(f, T)
  preimg = find(ismember(f, T));  % pre-imagem de T por f
end
\end{lstlisting}

% ---
\subsection{Transformações naturais entre funtores \texorpdfstring{$\Fin \to \Sub$}{Fin -> Sub}}

Dados dois funtores $F, G \colon \Fin \to \Sub$, uma transformação natural $\eta \colon F \Rightarrow G$ consiste numa família de morfismos em $\Sub$,
\[
  \eta_n \colon F(n) \hookrightarrow G(n) \quad \text{(i.e., } F(n) \subseteq G(n) \text{ para todo } n\text{)},
\]
tal que para todo morfismo $f \colon n \to m$ em $\Fin$:
\[
  G(f)\bigl(\eta_n(x)\bigr) \;=\; \eta_m\bigl(F(f)(x)\bigr), \quad \forall\, x \in F(n).
\]

\noindent\textbf{Exemplo:} Seja $F = P$ (funtor potência) e $G$ o funtor que associa a $n$ o conjunto de todos os multiconjuntos de $\{1,\ldots,n\}$. A inclusão natural de subconjuntos em multiconjuntos define uma transformação natural $\eta \colon P \Rightarrow G$.

% ============================================================
\section{Alínea (c) --- \texorpdfstring{$\Fin$}{Fin} e \texorpdfstring{$\Par$}{Par}}

\subsection{Descrição das categorias (modelo MATLAB/Octave)}

\begin{defbox}[Categoria $\Par$]
Um morfismo parcial $p \colon n \parto m$ é um vector de comprimento $n$ com entradas em $\{0, 1,\ldots,m\}$, onde $0$ significa que o elemento não tem imagem definida.
\end{defbox}

\begin{lstlisting}[language=Matlab, caption={Morfismo parcial em $\Par$}]
n = 4;  m = 3;
p = [2, 0, 1, 3];  % p(1)=2, p(2)=undef, p(3)=1, p(4)=3
\end{lstlisting}

% ---
\subsection{Composição em \texorpdfstring{$\Par$}{Par}}

A composição de duas aplicações parciais $p \colon n \parto m$ e $q \colon m \parto k$ é a aplicação parcial $(q \circ p) \colon n \parto k$ definida por
\[
  (q \circ p)(i) \;=\;
  \begin{cases}
    q(p(i)) & \text{se } p(i) \neq 0 \text{ e } q(p(i)) \neq 0,\\
    \text{indefinida} & \text{caso contrário.}
  \end{cases}
\]

\begin{lstlisting}[language=Matlab, caption={Composição de aplicações parciais}]
function r = compose_partial(p, q)
  r = zeros(1, length(p));
  for i = 1:length(p)
    if p(i) ~= 0 && q(p(i)) ~= 0
      r(i) = q(p(i));
    end
  end
end
\end{lstlisting}

% ---
\subsection{Funtor de inclusão \texorpdfstring{$I \colon \Fin \to \Par$}{I: Fin -> Par}}

Existe um funtor canónico de inclusão $I \colon \Fin \to \Par$ que envia cada conjunto finito $n$ a si próprio, e cada aplicação total $f \colon n \to m$ à mesma função vista como aplicação parcial (sem indefinições).

\begin{lstlisting}[language=Matlab, caption={Inclusão de $\Fin$ em $\Par$}]
function p = total_to_partial(f)
  p = f;  % aplicacao total e um caso particular de parcial
end
\end{lstlisting}

\noindent\textbf{Propriedades do funtor $I$:}
\begin{itemize}
  \item $I(\mathrm{id}_n) = \mathrm{id}_n$ em $\Par$ (vector $[1, 2, \ldots, n]$ sem indefinições).
  \item $I(g \circ f) = I(g) \circ I(f)$ (a composição parcial coincide com a total quando ambas são totais).
  \item $I$ é \emph{fiel}: morfismos distintos em $\Fin$ dão morfismos distintos em $\Par$.
  \item $I$ \emph{não é pleno}: há morfismos em $\Par$ que não provêm de $\Fin$ (os que têm indefinições).
\end{itemize}

% ---
\subsection{Transformações naturais \texorpdfstring{$\eta \colon F \Rightarrow G$}{eta: F => G} entre funtores \texorpdfstring{$\Fin \to \Par$}{Fin -> Par}}

Dados dois funtores $F, G \colon \Fin \to \Par$, uma transformação natural $\eta \colon F \Rightarrow G$ é uma família de morfismos em $\Par$,
\[
  \eta_n \colon F(n) \parto G(n), \quad \text{para cada objecto } n \in \Fin,
\]
satisfazendo a condição de naturalidade: para todo morfismo $f \colon n \to m$ em $\Fin$,
\[
  G(f) \circ \eta_n \;=\; \eta_m \circ F(f) \qquad \text{(como aplicações parciais)}.
\]

\noindent\textbf{Exemplo:} Seja $F = I$ o funtor de inclusão $\Fin \hookrightarrow \Par$, e $G$ o funtor que, dada uma aplicação total $f$, produz a mesma aplicação mas com o primeiro elemento do domínio tornado indefinido. Então $\eta_n \colon n \parto n$ pode ser a aplicação parcial $i \mapsto i$ para $i > 1$ e indefinida em $i = 1$. A naturalidade exige que esta perturbação seja compatível com todos os morfismos de $\Fin$.

% ============================================================
\section{Resumo Comparativo}

\begin{center}
\renewcommand{\arraystretch}{1.4}
\begin{tabular}{>{\bfseries}c p{3.2cm} p{5cm} p{4.5cm}}
\toprule
\rowcolor{darkblue!85}
\color{white}Alínea & \color{white}Par de categorias & \color{white}Funtor principal & \color{white}Transformação natural \\
\midrule
\rowcolor{lightblue!60}
(a) & $\Fin$, $\Vect$ & $F \colon \Fin \to \Vect$ (funtor livre) & $\eta_n \colon F(n) \to G(n)$, matrizes compatíveis \\
\rowcolor{white}
(b) & $\Fin$, $\Sub$ & $P \colon \Fin \to \Sub$ (cov.) \newline $P^* \colon \Fin^{\mathrm{op}} \to \Sub$ (contrav.) & $\eta_n \colon F(n) \subseteq G(n)$, inclusões naturais \\
\rowcolor{lightblue!60}
(c) & $\Fin$, $\Par$ & $I \colon \Fin \to \Par$ (inclusão) & $\eta_n \colon F(n) \parto G(n)$, aplic.\ parciais compatíveis \\
\bottomrule
\end{tabular}
\end{center}

\vfill
\begin{center}
\color{gray}\small Politécnico de Leiria --- MDPL26 Matemática Discreta, Pós-laboral
\end{center}

\end{document}
